\subsection{Observing the terminal metric potential}
\label{sec:metric-positive}
The spatial derivative of the remaining metric potential
measures an additional component of the orbit velocity.
This produces a positive result even at polarized
inhomogeneous backgrounds, where the two orbit transports
coincide.

For the linearized displacement of
\eqref{eq:transport-leading} put $X=\sqrt{c/a}\,V(c)$.
Linearizing the momentum constraint at $c$ gives
\[
 -\frac{\sqrt{c/a}}{2c}\,\delta\lambda_\theta(c)
  =\langle D_t X-X/(2c),U_\theta\rangle
   +\langle D_\theta X,U_t\rangle ,
\]
where $X$ in its time derivative denotes the full
time-rescaled displacement. The measured terminal orbit
fields determine the terms containing only $X$ and
its spatial derivative. Thus an observation of
$\lambda_\theta(c)$ gives $\langle D_tX,U_\theta\rangle$.
Its mean-zero spatial primitive has principal term
\[
 \tfrac12\left[\langle U_\theta,A_-\rangle H+
                    \langle U_\theta,A_+\rangle\mathcal T H\right].
\]
Indeed, the derivative of the two shifted primitives
in \eqref{eq:transport-leading} has the sum sign,
and commuting the spatial primitive through a smooth
coefficient gains an additional derivative.
The mean terms are finite rank.
The same forced-wave energy estimate used in
\eqref{eq:transport-leading} gives this principal
formula modulo a map from $H^{s-1}$ to $H^{s+1}$.

With $e_b$ the unit bisector from the mean rotation,
define $b_2=\cos\rho\,\langle U_\theta,e_b\rangle$.
Subtract the already observed translation difference
from the last principal expression. This controls
$b_2H$ in addition to $\beta H$.
Consequently the proof of
\eqref{eq:rank-two-compact} and
Theorem~\ref{th:fixed-rank-two} applies with
\[
 \inf_\theta\sum_{j=0}^{J}
  \big(|\beta(\theta-2jd)|^2+|b_2(\theta-2jd)|^2\big)>0
\]
and with $\cos\rho$ bounded away from zero.
For sufficiently short delays this holds if either
$P_\theta(a)$ or $Q_\theta(a)$ is nonzero somewhere.
For the first condition, $b_2$ converges uniformly to
$P_\theta(a)$ as $c\downarrow a$.
For the second, use \eqref{eq:relative-expansion}.
In either case the nonzero interval is covered by
finitely many translations. The resulting local estimate
uses the full terminal induced metric, plus a finite
low-mode $P$ reading if a lower-frequency kernel remains.

\begin{proposition}
\label{prop:two-metric-polarized}
Near a polarized background with nonconstant $P(a)$,
one sufficiently short rational terminal placement gives
a locally bi-Lipschitz missing-$P_t$ inverse from the
full terminal induced metric and the original initial
data. The same conclusion holds on an open class of
nearby Cauchy backgrounds. The observation norm includes
$\lambda_\theta(c)$ in $H^s$ in addition to
$P(c),Q(c)$ in $H^{s+1}$.
\end{proposition}

\begin{proof}
At the polarized reference, the $P$-variation solves
the scalar Euler--Poisson--Darboux equation and the $Q$
variation vanishes for the specified perturbation.
Apply the one-derivative scalar estimate at a short
rational gap as in Proposition~\ref{prop:safe-gap}
to bound the $H^{s-1}$ norm of the missing velocity by
the terminal $H^{s+1}$ norm of $P$.
Use it to remove the lower-order term of the
full-metric estimate above. Continuity of the
derivative and integration along a small neighborhood
give the local nonlinear statement.
For continuity on the unconstrained velocity space
one may use the periodic constraint source
$-2t\langle U_t,U_\theta\rangle$ as the output for all
reduced solutions. On periodic vacuum solutions it is
exactly $\lambda_\theta$. Restriction to initial
velocities compatible with the momentum constraint
therefore gives the claimed geometric observation.
\end{proof}

The analytic-path argument of Theorem~\ref{th:generic-two}
also applies to this full-metric output when $P(a)$ is
nonconstant. Keep $P(a)$ fixed, interpolate to a
polarized comparison state, and enforce the constraint
by subtracting a multiple of $P_\theta(a)$ from $P_t(a)$.
The geometric cover remains uniform along the path
because the same nonzero arc of $P_\theta(a)$ persists.
This gives nearby admissible backgrounds with stable
two-slice full-metric recovery at a common short delay
when either initial orbit field is spatially nonconstant.
